% ============================================================
%  Chapter 3 — Methodology
% ============================================================
\chapter{Methodology}
\label{ch:methodology}

\section{Study Area}
\label{sec:study_area}

The \studyarea{} occupies the northern part of Punjab province, a dissected
upland set between the Indus River to the west and the Jhelum River to the east.
For this study the plateau is defined by its five constituent administrative
districts---Chakwal, Rawalpindi, Attock, Jhelum, and Mianwali---which together
cover approximately \SI{25000}{\kilo\meter\squared} between roughly
\ang{32.16}--\ang{34.02}~N and \ang{71.11}--\ang{73.82}~E. The location and
extent of the area are shown in Figure~\ref{fig:intro_studyarea} of
Chapter~\ref{ch:introduction}.

Physiographically the plateau is far from uniform, and this heterogeneity is
central to the groundwater problem. The north, around Rawalpindi and Attock, is
higher, wetter, and underlain in places by coarse alluvial fans; the south and
southwest grade into the drier plains toward Mianwali and the folded sedimentary
ridges of the Salt Range and Kala Chitta Range. Drainage is organized chiefly
around the Soan River and its tributaries, the plateau's principal internal
watershed, together with the Haro and smaller systems
\cite{sarwar2022soandrought}. The climate is semi-arid, with mean annual
rainfall increasing broadly from the southwest toward the wetter northeast---on
the order of \SIrange{380}{600}{\milli\meter}, concentrated in the summer monsoon
between July and September \cite{hafeez2025rwh}. Because the area lies largely
outside the canal-irrigation command of the southern Punjab plains, agriculture
is predominantly rain-fed (\emph{barani}), and the population depends heavily on
groundwater drawn from both the alluvial aquifers along the river valleys and the
fractured hard-rock aquifers of the uplands \cite{waqas2022chakwal}. The
juxtaposition of permeable alluvium against near-impermeable crystalline and
folded rock is what makes groundwater potential vary so sharply across the
plateau, and is the spatial signal this study aims to capture.

\section{Methodological Overview}
\label{sec:overview}

The workflow proceeds in five stages, summarized in
Figure~\ref{fig:methodology}: (i)~collection of nine geo-environmental features
from satellite and reanalysis archives through Google Earth Engine (GEE);
(ii)~preprocessing of these layers into a co-registered, normalized feature
stack; (iii)~generation of three-class groundwater potential labels by a
knowledge-based weighted overlay; (iv)~training and comparison of four supervised
classifiers on a balanced sample, followed by application of the best model to
the full plateau; and (v)~validation of the resulting map against independent
evidence. Each stage is described in turn below.

\begin{figure}[H]
  \centering
  \includegraphics[width=0.92\textwidth]{Fig1_Methodology}
  \caption{Overall methodological workflow, from Google Earth Engine data
           collection through preprocessing, weighted-overlay label generation,
           machine-learning model training, prediction, and spatial-realism
           validation.}
  \label{fig:methodology}
\end{figure}

\section{Data Sources}
\label{sec:data_sources}

Nine features were selected to represent the terrain, climate, vegetation,
moisture, and soil controls on groundwater identified in
Chapter~\ref{ch:literature}. All were obtained through Google Earth Engine,
which provides analysis-ready access to the underlying archives and removes the
need to download and mosaic raw scenes. Table~\ref{tab:data_sources} lists each
feature, its source dataset, and its Pearson correlation with the groundwater
potential label (discussed further in Chapter~\ref{ch:results}).

\begin{table}[H]
\centering
\caption{The nine input features, their source datasets in Google Earth Engine,
         and their Pearson correlation ($r$) with the groundwater potential label.}
\label{tab:data_sources}
\small
\begin{tabular}{clll c}
\toprule
\# & Feature & Source & GEE dataset & $r$ \\
\midrule
1 & Elevation & SRTM            & \texttt{USGS/SRTMGL1\_003}                 & $+0.234$ \\
2 & Slope     & SRTM-derived    & \texttt{ee.Terrain.slope()}                & $+0.051$ \\
3 & Aspect    & SRTM-derived    & \texttt{ee.Terrain.aspect()}               & $+0.060$ \\
4 & TWI       & SRTM + HydroSHEDS & \texttt{WWF/HydroSHEDS/15ACC}            & $+0.185$ \\
5 & NDVI      & Landsat-8       & \texttt{LANDSAT/LC08/C02/T1\_L2}           & $+0.240$ \\
6 & NDWI      & Landsat-8       & \texttt{LANDSAT/LC08/C02/T1\_L2}           & $-0.222$ \\
7 & Rainfall  & CHIRPS          & \texttt{UCSB-CHG/CHIRPS/DAILY}             & $+0.514$ \\
8 & Soil      & SoilGrids       & \texttt{OpenLandMap/SOL/SOL\_TEXTURE}      & $+0.836$ \\
9 & LST       & MODIS           & \texttt{MODIS/061/MOD11A1}                 & $-0.426$ \\
\bottomrule
\end{tabular}
\end{table}

Three of these sources warrant specific mention. Soil texture was taken from
SoilGrids~2.0 \cite{poggio2021soilgrids}, a global digital soil-mapping product
generated with quantile regression forests trained on roughly 240,000 soil
profiles and more than 400 environmental covariates, with cross-validated,
spatially explicit uncertainty. It is the most comprehensive openly available
soil dataset, and---as Chapter~\ref{ch:results} shows---soil texture proved the
single strongest predictor in this study, consistent with its established control
on infiltration and aquifer permeability. The vegetation and water indices (NDVI
and NDWI) were derived from Landsat-8 Collection~2 surface-reflectance imagery
\cite{wulder2019landsat}. The terrain derivatives---elevation, slope, aspect, and
the topographic wetness index---were computed from the 30~m SRTM digital
elevation model; the choice of 30~m resolution is supported by Chang et al.
\cite{chang2019scaleeffects}, who found 30~m DEM derivatives to give the best
model accuracy among the resolutions they tested. The topographic wetness index
itself follows the standard TOPMODEL formulation of Beven and Kirkby
\cite{beven1979topmodel}, combining upslope contributing area (from the
HydroSHEDS flow-accumulation grid) with local slope, and cumulative rainfall was
summed from CHIRPS daily precipitation over the 2015--2023 period.

\section{Data Collection and Export}
\label{sec:collection}

Each feature was assembled in GEE as a single-band image, clipped to the union
of the five FAO/GAUL district polygons, and cast to floating point to avoid
type-mismatch errors when bands of different native types are combined. The nine
bands were stacked and exported as a single multi-band GeoTIFF at 30~m
resolution in geographic coordinates (EPSG:4326). The exported raster spans
$10{,}062 \times 6{,}884$ pixels---about 69.2~million cells---of which
\npixels{} (53.4\%) fall inside the district boundaries and carry valid data;
the remainder are nodata outside the study area. For metric calculations such as
area and distance, the data were reprojected to UTM Zone~42N (EPSG:32642).

\section{Preprocessing}
\label{sec:preprocessing}

Three preprocessing steps prepared the exported stack for analysis, all
implemented in Python with the \texttt{rasterio} library. First, pixels outside
the five district polygons were masked to nodata, fixing the valid-pixel count at
\npixels{}. Second, each band was independently rescaled to the interval
$[0,1]$ by min--max normalization (Equation~\ref{eq:minmax}), so that features
measured in very different units---metres, degrees, millimetres, index values---
contribute on a common scale and no single feature dominates by magnitude alone:

\begin{equation}
  x_{\text{norm}} = \frac{x - x_{\min}}{x_{\max} - x_{\min}}.
  \label{eq:minmax}
\end{equation}

Third, a quality check confirmed that all nine bands contained no missing values
within the valid mask and that every normalized band spanned the full $[0,1]$
range, ruling out constant or degenerate layers. The normalized stack and a
parallel table of pixel row/column indices (used later to reconstruct maps from
tabular predictions) formed the input to all subsequent steps.

\section{Feature Maps}
\label{sec:feature_maps}

Figure~\ref{fig:feature_panel} presents the nine normalized feature layers across
the plateau. Several of the spatial patterns anticipated in
Section~\ref{sec:study_area} are visible: rainfall and vegetation increase toward
the wetter north and northeast, surface temperature is highest over the exposed
southern terrain, and the wetness index picks out the valley floors and
low-lying tracts. These layers are the common input both to the weighted-overlay
labelling described next and to the machine-learning models.

\begin{figure}[H]
  \centering
  \includegraphics[width=\textwidth]{Fig3_Feature_Maps_Panel}
  \caption{The nine normalized geo-environmental input features across the
           \studyarea{}: elevation, slope, aspect, topographic wetness index,
           NDVI, NDWI, rainfall, soil texture, and land surface temperature.}
  \label{fig:feature_panel}
\end{figure}

\section{Label Generation by Weighted Overlay}
\label{sec:labels}

Because field measurements of groundwater are not available across the whole
plateau, reference labels for supervised training were generated using a
knowledge-based weighted overlay---an established approach in GIS-based spatial
mapping, and the same family of method used by the reference study
\cite{waqas2022chakwal} and validated as a labelling strategy by Roy and Saha
\cite{roy2019knowledgedriven} and others \cite{arulbalaji2019gisahp,pande2021ahpmif}.
Each normalized feature $f_i$ was assigned a weight $w_i$ reflecting its
hydrogeological influence on groundwater potential, and the weights were combined
into a single composite score $S$ for every pixel:

\begin{equation}
  S = \sum_{i=1}^{9} w_i \, f_i^{*},
  \label{eq:overlay}
\end{equation}

where $f_i^{*} = f_i$ for features that promote groundwater and
$f_i^{*} = 1 - f_i$ for features that suppress it. Features inversely related to
potential---slope, elevation, and land surface temperature---were inverted in
this way before weighting, so that, for instance, gentler slopes and cooler
(less evaporative) surfaces raise the score. The weights, listed in
Table~\ref{tab:weights}, give the largest influence to soil texture, the wetness
index, and rainfall, in line with both the hydrogeology of the area and the
correlations in Table~\ref{tab:data_sources}.

\begin{table}[H]
\centering
\caption{Feature weights and direction of influence used in the weighted-overlay
         label generation.}
\label{tab:weights}
\begin{tabular}{lcc}
\toprule
Feature & Weight (\%) & Direction \\
\midrule
Soil texture & 25 & Direct ($+$) \\
TWI          & 20 & Direct ($+$) \\
Rainfall     & 20 & Direct ($+$) \\
NDVI         & 10 & Direct ($+$) \\
Slope        & 10 & Inverse ($-$) \\
NDWI         &  5 & Direct ($+$) \\
Elevation    &  5 & Inverse ($-$) \\
LST          &  3 & Inverse ($-$) \\
Aspect       &  2 & Direct ($+$) \\
\midrule
\textbf{Total} & \textbf{100} & \\
\bottomrule
\end{tabular}
\end{table}

The continuous score was then divided into three ordinal classes using fixed
percentile thresholds: the lowest 40\% of scores were labelled \emph{Low}
potential (class~0), the 40th--80th percentile \emph{Medium} (class~1), and the
top 20\% \emph{High} (class~2). This 40/40/20 split reflects the expectation that
high-potential ground is the minority of a semi-arid plateau, and it fixes the
overall class balance of the reference map. The resulting label map is shown in
Figure~\ref{fig:label_map}.

\begin{figure}[H]
  \centering
  \includegraphics[width=0.85\textwidth]{gw_potential_label_map}
  \caption{Groundwater potential reference labels generated by the
           weighted-overlay model: Low (40\%), Medium (40\%), and High (20\%)
           potential zones across the \studyarea{}.}
  \label{fig:label_map}
\end{figure}

It is worth restating plainly what these labels are and are not. They encode
hydrogeological knowledge applied to the satellite features, not direct
observations of groundwater. A model trained on them therefore learns to
reproduce that knowledge-based product, and its internal accuracy measures
fidelity to the overlay rather than to ground truth. This is an accepted
limitation of GWPM in data-scarce regions; the response adopted here is not to
pretend otherwise but to validate the final map against independent evidence, as
described in Section~\ref{sec:validation_design} and Chapter~\ref{ch:results}.

\section{Training Dataset Construction}
\label{sec:training_data}

Training every model on all \npixels{} labelled pixels would be both
computationally wasteful and statistically skewed by the fixed 40/40/20 class
proportions. Instead, a balanced sample of 450,000 pixels was drawn---150,000
from each of the three classes---by stratified random sampling. Balancing the
classes prevents the models from favouring the majority classes and makes
metrics such as macro $F_1$ directly interpretable. The sample was partitioned
into 72\% training, 8\% validation, and 20\% test subsets, leaving 90,000 pixels
held out for final evaluation. The decision to fix and report this split
explicitly follows Nguyen et al. \cite{nguyen2021datasplitting}, who showed that
the partitioning scheme materially affects reported performance and should not be
left implicit.

\section{Machine-Learning Models}
\label{sec:ml_models}

Four classifiers were trained on the same data so that their performance could be
compared directly. They were chosen to span the two dominant tree-ensemble
strategies---bagging and boosting---and their combination:

\begin{itemize}
  \item \textbf{Random Forest (RF).} An ensemble of 300 decision trees grown on
        bootstrapped samples with Gini-impurity splits and a $\sqrt{9}$ feature
        subsample at each node \cite{liaw2007randomforest,tyralis2019randomforests}.
  \item \textbf{XGBoost.} Regularized gradient-boosted trees
        \cite{chen2016xgboost}, with learning rate 0.1, maximum depth 6, 300
        estimators, and 0.8 row/column subsampling.
  \item \textbf{LightGBM.} Histogram-based, leaf-wise gradient boosting
        \cite{ke2017lightgbm}, with 500 estimators, learning rate 0.05, and a
        minimum of 30 samples per leaf.
  \item \textbf{Voting Ensemble.} A soft-voting combination of the three models
        above with equal weights, averaging their predicted class probabilities.
\end{itemize}

All models were implemented in Python using \texttt{scikit-learn},
\texttt{xgboost}, and \texttt{lightgbm}. The comparison is deliberately aligned
with prior susceptibility-mapping work---most directly Sahin
\cite{sahin2020ensembletree}, who compared the same boosting and bagging
families---so that the model ranking obtained here can be read against an
external benchmark.

\section{Evaluation Metrics}
\label{sec:metrics}

Model performance was assessed on the 90,000-pixel test set using several
complementary measures rather than accuracy alone, since accuracy can be
misleading even on balanced data. Overall accuracy is reported alongside the
macro-averaged $F_1$-score, which weights all three classes equally; Cohen's
$\kappa$, which corrects agreement for chance; and the one-vs-rest ROC-AUC, which
summarizes class separability across decision thresholds. Per-class precision,
recall, and $F_1$ are also reported. The use of agreement- and
separability-based measures in addition to raw accuracy follows the
recommendation of Chicco and Jurman \cite{chicco2020mcc}. To confirm that the
results were not an artefact of a single partition, each model was additionally
evaluated under five-fold cross-validation on the training set.

\section{Prediction Map Generation}
\label{sec:prediction}

The best-performing model was then applied to all \npixels{} valid pixels to
produce the full-resolution groundwater potential map. Because holding the entire
feature stack in memory at once is impractical, prediction was carried out over
the raster in $2048 \times 2048$ pixel tiles, with each tile's pixels classified
in turn and written back to the output grid---twenty tiles in total, completing in
roughly twenty minutes. Two outputs were produced: a classified GeoTIFF assigning
each pixel to Low, Medium, or High potential, and a continuous raster giving the
model's predicted probability of the High class, which preserves the gradational
information lost in the hard classification and is used in the validation that
follows.

\section{Validation Design}
\label{sec:validation_design}

The central difficulty in validating this map is the one already noted: the
training labels are themselves model-derived, so a high agreement between the
prediction and the labels (reported in Chapter~\ref{ch:results}) confirms only
that the learning succeeded, not that the map is hydrogeologically correct. A
conventional well-by-well accuracy assessment was attempted but proved
uninformative, because the publicly available monitoring wells for the region are
not classified as productive or dry---they exist across both favourable and
unfavourable ground---and many fell on nodata pixels at the district margins.

The approach taken instead is what is termed here a \emph{spatial-realism
validation}: rather than asking whether the map matches a noisy point-by-point
ground truth, it asks whether the map's spatial structure agrees with independent
lines of hydrogeological evidence. Four such checks were designed:

\begin{enumerate}
  \item \textbf{Known productive sites.} Published groundwater sites in the
        Rawalpindi alluvial basin---ground that is known to be productive---were
        overlaid on the map, and the predicted class and High-probability at each
        were examined. A credible map should place these in the Medium or High
        classes.
  \item \textbf{Recharge-zone discrimination.} The independent network of
        rain-gauge stations compiled by Khan et al. \cite{khan2023rainfallrunoff}
        across the Pothohar basins was partitioned by hydrogeological setting into
        alluvial and non-alluvial groups, and the predicted High-probability was
        compared between groups using the Mann--Whitney $U$ test. Alluvial
        stations should show significantly higher predicted potential.
  \item \textbf{Cross-study consistency.} The predicted proportion of High
        potential was compared against the suitability proportions reported by an
        independent AHP study of the same region \cite{hafeez2025rwh}, to check
        that the overall magnitude is consistent with prior regional work.
  \item \textbf{District-scale gradient.} The fraction of each district classified
        as High potential was ranked and tested against the expected
        hydrogeological ordering (wetter, alluvial north versus drier, hard-rock
        south) using the Spearman rank correlation.
\end{enumerate}

Together these checks assess whether the map is right for the right reasons,
using evidence that played no part in its construction. Their results are
reported in Section~\ref{sec:validation_results}, and their interpretation, along
with the limitations of the approach, is taken up in
Chapter~\ref{ch:discussion}.
